\chapter*{Epilogue} \phantomsection \addcontentsline{toc}{chapter}{Epilogue. What Makes a Good Systematic Trader?} \markboth{Epilogue}{Epilogue}

TO FINISH THE BOOK I'D LIKE TO SUMMARISE WHAT I THINK ARE the main qualities a good systematic trader or investor should have.

A systematic trader should be humble, and underestimate their intelligence, skill and luck. Assume your trading will go badly, be prepared for that eventuality, and be pleasantly surprised if it doesn't. Don't try anything too clever; it is probably unnecessary and it's more likely to go wrong.

Use simple trading rules that have not been \textbf{over-fitted} or even fitted at all; constant no-rule forecasts, or discretionary forecasts combined with a strict stop loss policy. The \textbf{handcrafting} method of portfolio optimisation is both simple and effective.

You should be sceptical. Do not trust anybody. Don't trust the broker encouraging you to trade more, the trainer peddling you an expensive course, or the author\footnote{This includes me.} with a trading system that apparently worked for them, but might not be right for you.

Be pessimistic. Do not trust \textbf{back-tests}, even if you haven't over-fitted them, and even if they've been done on a \textbf{rolling out of sample} basis. The future is unlikely to be quite as good as the past. In my opinion a highly diversified system of systematic \textbf{trading rules} is unlikely to beat a \textbf{Sharpe ratio (SR)} of 1.0. If you're a \textbf{semi-automatic trader} who can't back-test their system, or an \textbf{asset allocating investor} with a \textbf{static} portfolio, then you should be even more pessimistic and assume maximum Sharpes of 0.5 and 0.4 respectively. You may do better, but you shouldn't expect to.

If you're making, or expect to make, steady gains with very few losses then there is a good chance you have a negative skew trading style which hasn't yet blown up, but probably will eventually. Adjust your \textbf{percentage} \textbf{volatility target} accordingly.

A good systematic investor will be thoughtful. You should know why you might be making money and why you might not. Understand your markets and your trading rules.

Thriftiness is another virtue. Know your trading costs. Stick to \textbf{speed limits}: I recommend you set your \textbf{turnover} so you never spend more than a third of the pessimistic expected Sharpe ratio for each \textbf{trading subsystem} on trading costs. Don't make your broker, or market makers, any wealthier.

You should be nervous. Only commit \textbf{trading capital} you can afford to lose. Use \textbf{Half-Kelly}: your maximum percentage volatility target should be half what you pessimistically expect your Sharpe ratio to be. At some point some instruments or trading rules in your portfolio will go badly wrong. Limit your exposure to this by trading the most diversified set of instruments and rules that you can. Stay away from low volatility instruments -- they're expensive to trade and dangerous.

The best systematic traders will be diligent when creating their systems, but lazy when running them. Put the hard work into designing a safe system that you are comfortable with and then do not change it. Make a commitment: don't be tempted to meddle, improve or risk manage. These time-consuming activities usually destroy performance.

Finally to make money you need to be lucky. Even if you do everything right you could still be unprofitable if the chance turns against you. You can't entirely eliminate risk from investing but you should quantify it, and make sure you can cope with the likely downside. So it only remains for me to say: Good Luck!
